\subsection{A falsifiable model of a veto}
\label{sec:context-design}
One possible cause of a reversal is a pair of tail features with opposing effects. Write $W(E)=1$ if restoration $E$ vetoes a fixed proposal. Suppose, for distinct $u,v\in T$, that its complete pattern is
\begin{equation}
 W(E)=\mathbf 1\{u\in E,\ v\notin E\}\quad\text{for every }E\subseteq T .
 \label{eq:pattern}
\end{equation}
This assumption has two observable consequences. Both endpoints must accept, and the same ordered pair must predict the outcome of every interior test. A single observed contradiction rules the pair out.

For a family $E_1,\dots,E_q$, let $C_u=\{j:u\in E_j\}$. The family observes a veto for orientation $(u,v)$ precisely when $C_u\setminus C_v$ is nonempty. Covering every orientation therefore requires pairwise incomparable signatures. Spencer's completely separating system attains the minimum number of informative queries, the smallest $q$ with $\binom{q}{\lfloor q/2\rfloor}\ge |T|$, by assigning distinct signatures of size $\lfloor q/2\rfloor$~\cite{spencer1970separating,hyttinen2013experiments}. We assign them in lexicographic order along the fitted tail ranking. Empty and full restorations are added as controls. Neither distinguishes two features. For five tail columns a single interior prefix of size two separates six of the 20 orientations. Four constant-weight probes separate all 20.

\paragraph{Can the actual performance metrics yield this pattern?}
For each panel take the largest excess of any metric's loss above its allowance in Eq.~\eqref{eq:ballot}. Let $g(E)$ be the second largest excess across the four panels. The vote vetoes exactly when $g(E)>0$. A sufficient condition for Eq.~\eqref{eq:pattern} is
\[
 g(E)=a+\delta(\mathbf 1_{u\in E}-\mathbf 1_{v\in E})+r(E),\quad
 |r(E)|\leq\eta,\quad a+\eta\leq0<a+\delta-\eta .
\]
This bound gives a positive excess only for the $u$-present, $v$-absent case. It depends on the units used for each metric. Our empirical test uses the weaker sign requirement alone.

For a binary example, take independent balanced signs $t_1,t_2,d,u,v$ and an independent uniform integer $z\in\{1,2,3,4\}$. The target is $Y=t_1t_2$. Keep $t_1,t_2,z$, delete $d$, and use zero as the masking value for the signs. Write $J_x=\mathbf 1\{x\ne0\}$ for visibility. A fixed classifier outputs $t_1t_2$ when $z\leq 2+J_d(J_u-J_v)$ and $-t_1t_2$ otherwise. On four equally weighted copies of the complete support, its symmetric class errors give identical accuracy and hard-score AUC. Both scores for the candidate are $1/2$. For the anchor they are $1/2+(J_u-J_v)/4$. Hence $g(E)=-0.015-10^{-12}+0.2425(J_u-J_v)$, which satisfies the sufficient condition with $\eta=0$. The classifier was deliberately made sensitive to the mask indicators.

\paragraph{What coverage misses.}
With equal-weight feature signatures, the number of missed ordered pairs is $\sum_c n_c(n_c-1)$ where $n_c$ columns share signature $c$. This count concerns Eq.~\eqref{eq:pattern} alone. For any other pattern a veto can occur on a single unqueried restoration. For instance, take a retained sign $t$, target $Y=t$, and a predictor equal to $t$ except that it adds one when $d$ has been masked and the tail equals a prescribed omitted set. Every queried MAE stays zero after deletion although the MAE on that set rises to one. Even perfect pair coverage can miss such a failure.
